% from ega1-5-fr.tex
\oldpage[I]{132}
Il y a en effet un ouvert $U$ de $Y$ tel que
$X=U\cap\overline{X}$~; comme $X$ est un sous-préschéma réduit de $Z$ en
vertu de (\hyperref[I.5.2.2-fr]{5.2.2}), le sous-préschéma $X$ est induit par
$Z$ sur le sous-espace ouvert $X$ en vertu de l'unicité
(\hyperref[I.5.2.1-fr]{5.2.1}).

\begin{corollary}[5.2.4]
\label{I.5.2.4-fr}
Soient $f:X\to Y$ un morphisme, $X'$ (resp. $Y'$) un sous-préschéma fermé de
$X$ (resp. $Y$) défini par un faisceau d'idéaux quasi-cohérent $\mathscr{J}$
(resp. $\mathscr{K}$) de $\mathscr{O}_X$ (resp. $\mathscr{O}_Y$). Supposons
que $X'$ soit réduit et que $f(X')\subset Y'$. Alors on a
$f^*(\mathscr{K})\mathscr{O}_X\subset\mathscr{J}$.
\end{corollary}

Comme la restriction de $f$ à $X'$ se factorise en
$X'\to Y'\to Y$ d'après (\hyperref[I.5.2.2-fr]{5.2.2}), il suffit d'appliquer
(\hyperref[I.4.4.6-fr]{4.4.6}).

\subsection{Diagonale; graphe d'un morphisme.}
\label{subsection:I.5.3-fr}

\begin{env}[5.3.1]
\label{I.5.3.1-fr}
Soit $X$ un $S$-préschéma~; on appelle \emph{morphisme diagonal} de $X$ dans
$X\times_S X$, et on note $\Delta_{X|S}$, ou $\Delta_X$, ou même $\Delta$ si
aucune confusion n'est possible, le $S$-morphisme $(1_X,1_X)_S$, autrement
dit, l'unique $S$-morphisme $\Delta_X$ tel que
\[
  p_1\circ\Delta_X=p_2\circ\Delta_X=1_X
  \tag{5.3.1.1}\label{I.5.3.1.1-fr}
\]
en désignant par $p_1$, $p_2$ les projections de $X\times_S X$
(déf. \hyperref[I.3.2.1-fr]{3.2.1}). Si $f:T\to X$, $g:T\to Y$ sont deux
$S$-morphismes, on vérifie aussitôt que
\[
  (f,g)_S=(f\times_S g)\circ\Delta_{T|S}
  \tag{5.3.1.2}\label{I.5.3.1.2-fr}
\]

Le lecteur observera que la définition précédente et les résultats énoncés
dans les n\textsuperscript{os} (\hyperref[I.5.3.1-fr]{5.3.1}) à (5.3.8) sont
valables dans toute catégorie, \emph{pourvu que les produits qui y figurent
existent dans cette catégorie}.
\end{env}

\begin{proposition}[5.3.2]
\label{I.5.3.2-fr}
Soient $X$, $Y$ deux $S$-préschémas~; si on identifie canoniquement le produit
$(X\times Y)\times(X\times Y)$ à $(X\times X)\times(Y\times Y)$, le morphisme
$\Delta_{X\times Y}$ s'identifie à $\Delta_X\times\Delta_Y$.
\end{proposition}

En effet, si $p_1$, $q_1$ sont les premières projections
$X\times X\to X$, $Y\times Y\to Y$, la première projection
$(X\times Y)\times(X\times Y)\to X\times Y$ s'identifie à $p_1\times q_1$,
et l'on a
\[
  (p_1\times q_1)\circ(\Delta_X\times\Delta_Y)
  =(p_1\circ\Delta_X)\times(q_1\circ\Delta_Y)=1_{X\times Y}
\]
même raisonnement pour les secondes projections.

\begin{corollary}[5.3.4]
\label{I.5.3.4-fr}
Pour toute extension $S'\to S$ du préschéma de base,
$\Delta_{X_{(S')}}$ s'identifie canoniquement à $(\Delta_X)_{(S')}$.
\end{corollary}

Il suffit de remarquer que $(X\times_S X)_{(S')}$ s'identifie canoniquement à
$X_{(S')}\times_{S'}X_{(S')}$
(\hyperref[I.3.3.10-fr]{3.3.10}).

\begin{proposition}[5.3.5]
\label{I.5.3.5-fr}
Soient $X$, $Y$ deux $S$-préschémas, $\varphi:S\to T$ un morphisme, faisant de
tout $S$-préschéma un $T$-préschéma. Soient $f:X\to S$, $Y\to S$ les
morphismes structuraux, $p$, $q$ les projections de $X\times_S Y$,
$\pi=f\circ p=g\circ q$ le morphisme structural $X\times_S Y\to S$. Alors, le
diagramme
\[
  \xymatrix{
    X\times_S Y\ar[r]^{(p,q)_T}\ar[d]_\pi &
    X\times_T Y\ar[d]^{f\times_T g}\\
    S\ar[r]_{\Delta_{S|T}} &
    S\times_T S
  }
  \tag{5.3.5.1}\label{I.5.3.5.1-fr}
\]
